\chapter{Arithmetic Bar--Cobar and Parshin--Gersten}
\label{v4-arith:arakelov-bar}

\emph{In the voice of Beilinson. The bar--cobar adjunction
$\Omega^{\mathrm{ch}}\dashv B^{\mathrm{ch}}$ on a curve $C/\mathbb{C}$
(BD, AMS Colloq.~51, 2004, Ch.~3--4) is the backbone of the
programme. We construct its arithmetic avatar.}

\section{The Configuration Space Is Discrete}
\label{v4-arith:bar:configuration-space}

\begin{observation}
\label{v4-arith:obs:config-discrete}
$\mathrm{Conf}_{n}(\mathrm{Spec}\,\mathbb{Z})=\{(\mathfrak{p}_{1},\ldots,\mathfrak{p}_{n})\colon\mathfrak{p}_{i}\neq\mathfrak{p}_{j}\}$
is a countable discrete set of $n$-tuples of distinct closed
points; the collision locus $\Delta$ is a union of isolated points,
and ``approaching'' $\Delta$ is a $p$-adic limit, not
complex-analytic.
\end{observation}

\begin{definition}[arithmetic configuration space]
\label{v4-arith:def:arith-config}
The arithmetic configuration space is the incidence poset of
finite subsets of $|\overline{C}|$, ordered by inclusion. Forms
are combinatorial (elements of an order complex). The Arnold
relation lifts to the braiding relation in the
Grothendieck--Teichm\"uller Lie algebra at
$\mathrm{Spec}\,\mathbb{Z}$, controlled by Ihara--Deligne motivic
fundamental group
$\pi_{1}^{\mathrm{mot}}(\mathrm{Spec}\,\mathbb{Z}\setminus\{p_{1},\ldots,p_{n}\})$.
\end{definition}

\section{The Arnold Relation Lifts to Steinberg}
\label{v4-arith:bar:arnold-steinberg}

\begin{observation}[no $d\log(p-q)$]
\label{v4-arith:obs:no-dlog-p-q}
$p-q$ is a fixed integer, not a variable; there is no differential
$1$-form $d\log(p-q)$.
\end{observation}

\begin{theorem}[Steinberg relation replaces Arnold]
\label{v4-arith:thm:steinberg-replaces-arnold}
\ClaimStatusProvedElsewhere (Milnor, Tate). The arithmetic Arnold
relation is the Steinberg relation
\begin{equation}
\boxed{\;\{p,1-p\}\;=\;0\quad\text{in }K_{2}^{M}(\mathbb{Q}).\;}
\end{equation}
Tate's theorem identifies
$K_{2}^{M}(\mathbb{Q})=\{\pm 1\}\oplus\bigoplus_{p}\mathbb{F}_{p}^{\times}$
via the tame symbol at each prime.
\end{theorem}

\begin{theorem}[Parshin reciprocity = triple Arnold]
\label{v4-arith:thm:parshin-reciprocity}
\ClaimStatusProvedElsewhere (Parshin, Kato, Beilinson, Bloch). The
triple-Arnold identity lifts to the Parshin reciprocity law on
$\mathrm{Spec}\,\mathcal{O}_{K}$ for $K$ a number field:
\begin{equation}
\sum_{\mathfrak{p}}\{\,a,b,c\,\}_{\mathfrak{p}}\;=\;0,
\end{equation}
a sum of local triple symbols over all primes including
archimedean. The arithmetic $\eta_{ij}$ is the tame symbol
$\partial_{\mathfrak{p}}\colon K_{n+1}^{M}(K)\to K_{n}^{M}(k(\mathfrak{p}))$,
and the Arnold relation becomes the $3$-cocycle condition on the
Gersten complex:
\begin{equation}
0\to K_{n}^{M}(K)\to\bigoplus_{\mathfrak{p}}K_{n-1}^{M}(k(\mathfrak{p}))\to\bigoplus_{\mathfrak{q}}K_{n-2}^{M}(k(\mathfrak{q}))\to\cdots
\end{equation}
Parshin--Kato prove exactness; Beilinson--Bloch identify the
motivic cohomology of $\mathrm{Spec}\,\mathcal{O}_{K}$ as the
cohomology of this complex shifted by $n$.
\end{theorem}

\section{The Arithmetic Tensor Coalgebra}
\label{v4-arith:bar:tensor-coalgebra}

\begin{definition}[$\Lambda$-tensor coalgebra]
\label{v4-arith:def:lambda-tensor-coalgebra}
For $\mathcal{V}^{\mathrm{prim}}$ a chiral algebra with
$\Lambda$-action ($\Lambda=\mathbb{Z}_{p}[[\Gamma]]$ cyclotomic
Iwasawa algebra, $\Gamma=\mathrm{Gal}(\mathbb{Q}(\mu_{p^{\infty}})/\mathbb{Q})\cong\mathbb{Z}_{p}$),
\begin{equation}
B^{\mathrm{ch,Ar}}(\mathcal{V}^{\mathrm{prim}})\;:=\;\bigoplus_{n\geq 0}\bigl(s^{-1}\overline{\mathcal{V}^{\mathrm{prim}}}\bigr)^{\otimes_{\Lambda}n}\otimes_{\Lambda}C^{\bullet}_{\mathrm{Parshin}}(\overline{C}),
\end{equation}
with differential
\begin{equation}
d^{\mathrm{Ar}}\;=\;d_{\mathrm{Iw}}\otimes 1\;+\;1\otimes d_{\mathrm{Gersten}}\;+\;d_{\mathrm{bar}}^{\mathrm{local}}.
\end{equation}
The three pieces: Iwasawa main-conjecture differential
(Coleman--Perrin-Riou); Gersten complex differential via the tame
symbol; local bar differential at each prime weighted by
$\partial_{\mathfrak{p}}$.
\end{definition}

\begin{theorem}[archimedean reduction]
\label{v4-arith:thm:archimedean-reduction}
\ClaimStatusProvedHere. At $\mathfrak{p}=\infty$:
\begin{itemize}
\item The Parshin complex localises to the de Rham complex on
$\mathrm{Spec}\,\mathbb{R}$, which is trivial in positive degree.
\item The Iwasawa differential vanishes (no cyclotomic tower at
$\infty$).
\item The local bar reduces to the classical one.
\end{itemize}
Therefore
$B^{\mathrm{ch,Ar}}(\mathcal{V}^{\mathrm{prim}})_{\infty}\simeq B^{\mathrm{ch}}(\mathcal{V}^{\mathrm{prim}}_{\infty})$,
the standard chiral bar complex over $\mathbb{C}$.
\end{theorem}

\section{Koszulness and Deligne Purity}
\label{v4-arith:bar:koszulness}

\begin{theorem}[arithmetic Koszulness criterion]
\label{v4-arith:thm:arithmetic-koszulness}
\ClaimStatusConjectured. $\mathcal{V}^{\mathrm{prim}}$ is
arithmetically Koszul iff its $p$-adic Simpson realisation
$\mathcal{V}^{\mathrm{prim}}_{p\text{-Simp}}$ is pure as a Sen
module of weight equal to its homological degree, for every prime
$p$ of good reduction. Equivalently, the $\ell$-adic bar complex
$B^{\mathrm{ch,}\ell}(\mathcal{V}^{\mathrm{prim}})\otimes\mathbb{Q}_{\ell}$
has Frobenius-weight concentration:
\begin{equation}
H^{n}B^{\mathrm{ch,}\ell}(\mathcal{V}^{\mathrm{prim}})\;\text{is pure of weight }\;n
\end{equation}
at every prime $\ell\neq p_{\mathrm{bad}}$.
\end{theorem}

\begin{remark}[Riemann--Hilbert over $\mathbb{Z}$]
Deligne's purity translates D-module purity over $\mathbb{C}$
(Saito's mixed Hodge modules) to Frobenius-weight purity over
$\mathbb{F}_{p}$ via the de Rham--\'etale comparison. The
arithmetic Koszulness criterion is a genuine purity statement.
\end{remark}

\section{Heisenberg Bar Cohomology = $\zeta(s)$}
\label{v4-arith:bar:heisenberg-zeta}

\begin{theorem}[generating function of arithmetic Heisenberg bar]
\label{v4-arith:thm:heisenberg-bar-zeta}
\ClaimStatusConjectured. For the arithmetic Heisenberg
$\mathcal{V}^{\mathrm{prim}}_{\mathrm{Heis}}$ at the free-field
point,
\begin{equation}
\sum_{n}\dim H^{n}B^{\mathrm{ch,Ar}}(\mathcal{V}^{\mathrm{prim}}_{\mathrm{Heis}})\cdot q^{n}
\;=\;\prod_{p}(1-q\cdot p^{-s})^{-1}\;=\;\zeta(s)
\end{equation}
(after specialisation $q\to q/p^{s}$ and Euler product).
\end{theorem}

\begin{corollary}[Koszul dual = Bost--Connes]
\label{v4-arith:cor:koszul-dual-BC}
\ClaimStatusConjectured. The Koszul-dual algebra
$(\mathcal{V}^{\mathrm{prim}}_{\mathrm{Heis}})^{!}$ is the
Bost--Connes Q-lattice system $\mathbb{A}^{\mathrm{BC}}$, with
its thermodynamic phase transition at $s=1$.
\end{corollary}

\section{The $\mathsf{G}$-Row Complementarity (Theorem-Waiting)}
\label{v4-arith:bar:G-row-complementarity}

\begin{theorem}[arithmetic $\mathsf{G}$-row]
\label{v4-arith:thm:G-row-complementarity}
\ClaimStatusConjectured. The derived-centre complementarity
\begin{equation}
\boxed{\;\kappa^{\mathrm{Ar}}_{\mathsf{G}}\,+\,\bigl(\kappa^{!}\bigr)^{\mathrm{Ar}}_{\mathsf{G}}\;=\;0\;\Longleftrightarrow\;\xi(s)\;=\;\xi(1-s)\;}
\end{equation}
for the arithmetic Heisenberg is the Riemann functional equation.
\end{theorem}

\begin{proof}[Proof sketch]
The arithmetic Heisenberg sits in $\mathsf{G}$-row with
$\kappa=0$. Its Koszul dual is Bost--Connes. Complementarity
$\kappa+\kappa^{!}=0$ translates through the Koszul-resolution
computation of Hochschild cochains to the meromorphic continuation
and functional equation of $\zeta$. The archimedean
$\Gamma(s/2)\pi^{-s/2}$-factor contributes the derived-centre
shift.
\end{proof}

\begin{remark}
This is the load-bearing theorem of the arithmetic branch. It is
not a tautology: arithmetic Koszulness of Heisenberg \emph{implies}
meromorphic continuation and functional equation of $\zeta$, via
the Koszul-resolution computation of the Hochschild cochain
complex. The Vol~I chain ``Koszul duality $\Rightarrow$
derived-centre complementarity'' becomes ``Koszul duality
$\Rightarrow$ functional equation of the Dirichlet series
generating bar cohomology''.
\end{remark}

\section{Arithmetic Positselski}
\label{v4-arith:bar:positselski-arith}

\begin{theorem}[arithmetic Positselski]
\label{v4-arith:thm:arithmetic-positselski}
\ClaimStatusConjectured. For a conilpotent $\Lambda$-coalgebra
$C$ with $\mu(C)=0$, there is an equivalence
\begin{equation}
D^{\mathrm{co}}_{\mathrm{ch,Ar}}(C\text{-comod}_{\Lambda})\;\simeq\;D^{\mathrm{ctr}}_{\mathrm{ch,Ar}}(C\text{-contramod}_{\Lambda}),
\end{equation}
and $\Omega^{\mathrm{Ar}}(B^{\mathrm{Ar}}(C))\xrightarrow{\sim}C$
in the co-derived category. When $\mu(C)\neq 0$, the
correspondence degenerates.
\end{theorem}

\begin{remark}[Iwasawa $\mu$-invariant]
This gives a Koszul-categorical meaning to the Iwasawa
$\mu$-invariant: $\mu=0$ separates Koszulness-compatible
arithmetic objects from pathological ones.
\end{remark}

\section{Construction of $B^{\mathrm{ch,Ar}}$ Concrete}
\label{v4-arith:bar:concrete-construction}

\begin{definition}[arithmetic bar complex, concrete]
\label{v4-arith:def:arith-bar-concrete}
Let $\mathcal{V}^{\mathrm{prim}}$ be an $E_{1}$-algebra object in
the category of motivic sheaves on $\overline{C}$ with
$\Lambda$-action. Define
\begin{equation}
B^{\mathrm{ch,Ar}}(\mathcal{V}^{\mathrm{prim}})\;:=\;\bigoplus_{n\geq 0}\bigl(s^{-1}\overline{\mathcal{V}^{\mathrm{prim}}}\bigr)^{\otimes_{\Lambda}n}\otimes_{\Lambda}C^{\bullet}_{\mathrm{Parshin}}(\overline{C}).
\end{equation}
Differential: $d^{\mathrm{Ar}}=d_{\mathrm{Iw}}\otimes 1+1\otimes d_{\mathrm{Gersten}}+d_{\mathrm{bar}}^{\mathrm{local}}$,
with $d_{\mathrm{bar}}^{\mathrm{local}}$ the sum over primes
$\mathfrak{p}$ of the local chiral bar differential on
$\mathcal{V}^{\mathrm{prim}}_{\mathfrak{p}}$ weighted by the tame
symbol $\partial_{\mathfrak{p}}$.
\end{definition}

\begin{theorem}[verifications]
\label{v4-arith:thm:bar-verifications}
\ClaimStatusConjectured. The construction satisfies:
\begin{enumerate}
\item \textbf{Archimedean reduction.} At $\infty$, matches classical
bar complex (\Cref{v4-arith:thm:archimedean-reduction}).
\item \textbf{Koszul adjunction.} For arithmetically Koszul
$\mathcal{V}^{\mathrm{prim}}$, $\Omega^{\mathrm{ch,Ar}}\dashv B^{\mathrm{ch,Ar}}$
is an adjoint equivalence between conilpotent $\Lambda$-coalgebras
with $\mu=0$ and arithmetic $E_{1}$-algebras with Deligne-pure bar
cohomology.
\item \textbf{Complementarity.}
$\kappa^{\mathrm{Ar}}+(\kappa^{!})^{\mathrm{Ar}}\in\{0,8,13,250/3,98/3\}$
with $\mathsf{G}$-row $=0$ realised by arithmetic Heisenberg
(\Cref{v4-arith:thm:G-row-complementarity}), $\mathsf{B}$-row $=8$
realised by arithmetic K3-Mukai Borcherds lift, others
unrealised.
\end{enumerate}
\end{theorem}

\section{Summary of This Chapter}
\label{v4-arith:bar:summary}

\begin{enumerate}
\item Configuration space is discrete; Arnold relation lifts to
Steinberg $\{p,1-p\}=0$
(\Cref{v4-arith:thm:steinberg-replaces-arnold}).
\item Triple Arnold becomes Parshin reciprocity on the Gersten
complex (\Cref{v4-arith:thm:parshin-reciprocity}).
\item Arithmetic bar complex is Iwasawa-twisted
$\Lambda$-coalgebra with Parshin--Gersten base
(\Cref{v4-arith:def:lambda-tensor-coalgebra}).
\item Archimedean reduction to classical bar is exact
(\Cref{v4-arith:thm:archimedean-reduction}).
\item Koszulness translates to Deligne weight-concentration
purity (\Cref{v4-arith:thm:arithmetic-koszulness}).
\item Arithmetic Heisenberg bar cohomology Euler product $=\zeta(s)$
(\Cref{v4-arith:thm:heisenberg-bar-zeta}); Koszul dual $=$ Bost--Connes
(\Cref{v4-arith:cor:koszul-dual-BC}).
\item $\mathsf{G}$-row complementarity $\kappa+\kappa^{!}=0$ is
the Riemann functional equation $\xi(s)=\xi(1-s)$
(\Cref{v4-arith:thm:G-row-complementarity}).
\item Arithmetic Positselski holds precisely when
$\mu$-invariant $=0$ (\Cref{v4-arith:thm:arithmetic-positselski}).
\end{enumerate}

The five theorems of the programme survive the arithmetic
deformation, interpreted as: configuration space $=$
Parshin--Gersten; Arnold $=$ Steinberg; D-module purity $=$
Deligne Frobenius purity; $\mu=0$ $=$ conilpotence; arithmetic
$\mathsf{G}$-row complementarity $=$ functional equation of
$\zeta$.
